\begin{longtable}{@{}p{0.4cm}p{3.1cm}p{5.4cm}p{6.2cm}@{}}
\caption{Summary of the literature most relevant to this framework.}\label{tab:lit}\\
\toprule
\textbf{\#} & \textbf{Reference} & \textbf{Focus / method} & \textbf{What we reuse} \\
\midrule\endfirsthead
\multicolumn{4}{@{}l}{\small\itshape Table \ref{tab:lit}, continued}\\
\toprule
\textbf{\#} & \textbf{Reference} & \textbf{Focus / method} & \textbf{What we reuse} \\
\midrule\endhead
\bottomrule\endlastfoot
1 & \small Rashida \emph{et al.} (2021), \emph{Computers} 10(5):57 & \small A multi-method framework for evaluating university websites in Bangladesh. The automated tool crawls with Selenium and evaluates three criteria: content of information (25 key strings), website performance and loading time (20 browser timing variables). A questionnaire study compares the survey ranking with the automated one. & \small The principal baseline. Most of our content attributes --- vision and mission, faculty, notice board, library, contact --- originate in or extend its content criteria. We extend its automated evaluation from a national to a global corpus, and replace its fixed weighting with a documented, sensitivity-tested one. \\\addlinespace[3pt]
2 & \small Saleh \emph{et al.} (2022), \emph{Indonesian J.\ of Electrical Engineering and Computer Science} 25(1):511--520 & \small A systematic literature review of university website quality evaluation. It extracts 79 quality factors from prior work and consolidates them into six: information quality, specific content, usability, web appearance, service interaction quality and functionality. It reports the absence of any universally accepted model. & \small The six-factor structure is the theoretical foundation for grouping and justifying our attribute blocks. Its finding that no accepted model exists is what obliges us to construct a label and to state its construction in full (Section~\ref{sec:label}). \\\addlinespace[3pt]
3 & \small Mukanda, Mbuguah \& Wabwoba (2022), \emph{Int.\ J.\ of Computer Trends and Technology} 70(2):1--9 & \small Evaluates the usability of the top five Kenyan university websites using automated tools (Nibbler, Silktide, Semrush, Ahrefs), covering responsiveness, colour contrast, broken links, speed, compression, SEO, and accessibility. & \small The source of our technical and accessibility attributes: mobile responsiveness, contrast ratio, broken links, load speed, compression, SEO metadata and alt-text coverage. Where it evaluates five universities, we evaluate 1,225. \\\addlinespace[3pt]
4 & \small Biyyapu \emph{et al.} (2023), \emph{Computers} 12(9):181 & \small Applies machine learning to website quality classification from automatically extracted completeness attributes (missing URLs, images, videos, PDFs, form fields), comparing an MLP against SVM and other classifiers. & \small Motivates the move from a fixed rule-based score to a learned model, and the practice of comparing several algorithms on one feature vector. Our feature space is substantially broader, and Section~\ref{sec:why} goes beyond reporting which family wins to identifying the structural reason it wins. \\\addlinespace[3pt]
\end{longtable}
